\documentclass[11pt,a4paper]{article}
\usepackage[utf8]{inputenc}
\usepackage{amsmath,amssymb,amsfonts}
\usepackage{graphicx}
\usepackage{hyperref}
\usepackage{booktabs}
\usepackage{cite}
\usepackage{microtype}
\usepackage[margin=1in]{geometry}

\title{Kvant Tizimlarida Chuqur Neyron Tarmoqlari}
\author{O'zbekiston Milliy Tadqiqot Laboratoriyasi}
\date{\today}

\begin{document}
\maketitle

% --- Section: 00_abstract.tex ---
\section{Abstract}

We propose an empirical evaluation of attention mechanisms in quantum simulations.

% --- Section: 01_introduction.tex ---
\section{Introduction}

Deep learning models have accelerated computational physics.

% --- Section: 02_methodology.tex ---
\section{Methodology}

We consider the Hamiltonian formulation E = mc^2.

% --- Section: 03_results.tex ---
\section{Results}

of attention mechanisms in quantum simulations.
1. Introduction
Deep learning models have accelerated computational physics.
2. Methodology
We consider the Hamiltonian formulation E = mc^2.
3. Results
Empirical accuracy reached 98\% on benchmark datasets.

% --- Section: 04_conclusion.tex ---
\section{Conclusion}

Our architecture provides strong theoretical guarantees.

\bibliographystyle{plain}
\bibliography{references}
\end{document}